\scanpage{23}
% Pencil foliation: 22, lower left.
\uline{Motifs}

1) À tout préschéma noeth (éventuellement de type fini sur un anneau
noethérien) $X$ est associé une \uline{catégorie abélienne}
$\mathcal{M}^+(X)$, dite catégorie des \uline{motifs} \inserted{\uline{effectifs}} sur $X$.
C'est une
$\mathbb{Q}$-catégorie \deleted{de $\mathbb{Q}$-modules} divisible, i.e. pour tt
$M\in\operatorname{Ob}\mathcal{M}^+(X)$ et $n\in\mathbb{Z}$, $n\ne0$,
$n\,\operatorname{id}_M$ est un isom
de $M$. De plus, $\mathcal{M}^+(X)$ est munie d'un
produit tensoriel commutatif et unitaire,
\inserted{(exact à droite,)}
l'unité est notée $\mathbf1_X$ ou $\mathbb{Q}_X(0)$

\margin{On peut en
termes des
\unsure{données} construire
des foncteurs
$\lambda^iM$ etc \dots}

On considère aussi la catégorie dérivée bornée
$D^b(\mathcal{M}^+(X))$ de $\mathcal{M}^+(X)$. Le produit tensoriel est
étendu en un bifoncteur $M$\uuline{$\otimes$}$N$ en $M,N\in D^b(\mathcal{M}^+(X))$

2) Soit $f:X\longrightarrow Y$ un morphisme de préschémas
\deleted{\unsure{loc}} noeth, il lui est associé un \uline{foncteur}
\uline{exact}
\diagram{\begin{tikzcd}
f^*:\mathcal{M}^+(Y)\arrow[r]&\mathcal{M}^+(X)
\end{tikzcd}}
compatible avec $\otimes$,
\inserted{d'où
\diagram{\begin{tikzcd}
Lf^*:D^b(\mathcal{M}^+(Y))\arrow[r]&D^b(\mathcal{M}^+(X)),
\end{tikzcd}}}
On a transitivité.
Si $f$ est de type fini, et propre \uline{ou} $Y$ excellent,
on \deleted{définit} a de \unsure{$=$|$\overline{\mathrm m}$} un foncteur
\diagram{\begin{tikzcd}
Rf_*:D^b(\mathcal{M}^+(X))\arrow[r]&D^b(\mathcal{M}^+(Y))
\end{tikzcd}}
